% =============================================================================
% 档案员按语 — Diegetic Frame of Archive 1,260
% =============================================================================
\cleardoublepage
\thispagestyle{empty}
\chapter*{档案员按语}
\addcontentsline{toc}{chapter}{档案员按语}
\markboth{档案员按语}{}

\begin{center}
\itshape\small
此信封内装有患者1,260日记的完整誊本。\\
阅读原始材料在检验人员身上引起了心身改变。\\[0.5em]
\upshape\scriptsize --- 法证信封第1,260/1979号上的钢笔批注
\end{center}

\vspace{1.2em}

\begingroup

您手中这部书并非为出版而写。它是巴巴塞纳收容院患者1,260\textbf{若阿金·安东尼奥·科雷亚}日记的法证重建。

1915年大旱与塞阿拉难民集中营的幸存者若阿金，于1917年被送入收容院，并在那里度过了62年，直到1979年去世。在那段监禁岁月中，他写下了847页零散笔记。遗体从乱葬坑中被掘出时——下葬62天后——原笔记本用布条绑在他的胸前。

\vspace{0.4em}
\noindent 根据掘墓人的法证陈述，纸张散发出潮湿与泥土的甜味，墨迹似乎永远不会干。验尸报告将该发现登记为保管物品，并附入病历。

\vspace{0.6em}
本版汇集了按五个维度组织的78个碎片：塞尔唐童年的记忆；精神病院的隔离与折磨；临床报告的检视；卢西亚·门德斯医生于2026年进行的法证调查；以及读者自身受污染的记录。

\vspace{0.6em}
这不是一个简单的恐怖故事。这些书页中记载的痛苦不是虚构：它跨越了数十年，消耗了医生和检验人员，并在阅读这一行为中寻找新的宿主。

\vspace{1.2em}
\noindent{\itshape
建议不要在黑暗中阅读本书。\\[0.35em]
建议不要完全独处时阅读本书。\\[0.35em]
若阅读中升起湿土的气味，请立即合上本书。}

\vspace{1.2em}
\noindent 这段叙事所依托的历史创伤是真实的。巴巴塞纳收容院（1903—1994）未经同意收容了六万多人；其乱葬坑中掘出1,853具遗体。若阿金是一个角色，但他所承载的痛苦——饥饿、流离、囚禁、精神科酷刑——属于成千上万无名的巴西人。

\vspace{0.6em}
\noindent 此处唤起的恐怖不应作为娱乐消费。对读者的操纵是一种伦理与形式的选择——不是为了欺骗，而是为了让身体感受到心灵已知之事：乱葬坑不是隐喻，被记载的苦难存在是为了被见证，而不是被审美化。

\vspace{0.6em}
\noindent{\itshape 您将被操纵。您将被污染。但到头来，您会知道为什么。}

\vspace{1.4em}
\hfill{\scshape --- 档案员}\hspace{1em}

\endgroup
\clearpage
